\section{Significance of Research}
% =============================================================================
% SIGNIFICANCE - v3 EMERGENT SOCIAL STRUCTURES FRAMING
% =============================================================================
% STATUS: Needs update to reflect broader framing
%
% KEY POSITIONING: Extending Axelrod into the semantic reasoning era
%
% SIGNIFICANCE AREAS:
%
% 1. EXTENDING CLASSIC COOPERATION RESEARCH
%    - Axelrod (1984) showed optimizers can cooperate mechanistically
%    - You extend: do semantic reasoners cooperate DIFFERENTLY?
%    - This is a natural evolution of foundational work
%    - Positions you in established conversation, not starting from scratch
%
% 2. MULTI-AGENT AI SYSTEMS
%    - LLM agents increasingly deployed in multi-agent settings
%    - Need to understand: do they coordinate differently than RL?
%    - Practical implications for system design
%    - AI safety relevance: can semantic reasoning enable better coordination?
%
% 3. METHODOLOGICAL CONTRIBUTION
%    - Two-comparison design isolates architecture vs framing
%    - Addresses confounding problem in RL vs LLM comparisons
%    - Provides replicable framework for future studies
%    - Expanded metrics for emergent structure analysis
%
% 4. THEORETICAL IMPLICATIONS (broader than IR)
%    - Tests whether semantic priors enable different social structures
%    - Relevant to: organizational behavior, market dynamics, social simulation
%    - IR (neorealism vs constructivism) is ONE application domain
%
% 5. TIMELINESS
%    - LLMs are new; systematic comparisons to RL needed
%    - Growing use of LLM agents in research and deployment
%    - Understanding their emergent behavior is urgent
%
% EXAMPLE DIRECTION:
%
%   "This research extends the classic study of cooperation among self-interested
%   agents (Axelrod 1984) into the era of semantic reasoning. While Axelrod
%   demonstrated that optimizers can discover cooperation mechanistically, we
%   ask whether semantic priors enable qualitatively different social structures.
%
%   The contribution is threefold:
%
%   First, for multi-agent AI systems, understanding whether LLM agents coordinate
%   differently than RL agents has practical implications. As LLMs are deployed
%   in increasingly complex multi-agent settings, knowing whether semantic priors
%   produce different emergent structures matters for system design and safety.
%
%   Second, methodologically, the two-comparison design provides a framework for
%   isolating reasoning architecture effects from goal framing effects. This
%   addresses a common confound in RL vs LLM comparisons.
%
%   Third, theoretically, this work tests whether reasoning processes shape
%   emergent social structures - a question relevant beyond any single domain.
%   The results inform debates in international relations, organizational behavior,
%   and computational social science about whether ideational factors can overcome
%   material incentive structures."
%
% TYPO TO FIX: "comparitive" → "comparative"
% =============================================================================

With the debate between Neorealism and Constructivism still actively ongoing, my research could potentially provide additional insights. Though, ABM models are widely used in IR studies, the Artificial Intelligence approach to designing the agent architecture seems quite novel. Especially, comparitive studies between RL and LLM behaviour in social studies is an active field of research.

% TODO: Rewrite for v3 emergent structures framing:
%
% STRUCTURE:
%   1. Position as extending Axelrod (established, foundational)
%   2. Multi-agent AI relevance (practical, timely)
%   3. Methodological contribution (two-comparison design)
%   4. Theoretical implications (broader than IR)
%   5. Timeliness (LLMs are new, need systematic understanding)
%
% KEY FRAMING:
%   - Not "contributing to IR debate" (narrow)
%   - But "extending cooperation research into semantic reasoning era" (broad)
%   - IR is ONE application domain, not the main frame
%
% STRONGER CLAIMS:
%   - "First systematic comparison of emergent structures from RL vs LLM
%      agents with controlled experimental design"
%   - "Extends Axelrod's foundational work on cooperation to semantic reasoning"
%   - "Provides replicable framework for studying reasoning architecture effects"
